\section{原始関数と不定積分}\label{節:原始関数基本定理}
基本定理により,積分で作った関数 $A(x)$ を微分すると $f$ に戻ることが分かった.では逆に,微分すると $f$ になる関数を先に1つ見つけておけば,定積分を計算できるのではないだろうか.

\begin{definition}{原始関数と不定積分}
1つの区間で $F'(x)=f(x)$ を満たす関数 $F$ を,\,$f$ の\textbf{原始関数}という.すべての原始関数を表す記法として
\[
 \int f(x)\dd x=F(x)+C
\]
と書き,\textbf{不定積分}と呼ぶ.\,$C$ は任意の定数であり,\textbf{積分定数}という.
\end{definition}

なぜ定数を足すだけで,すべての原始関数を表せるのだろうか.\,$F,G$ がともに $f$ の原始関数なら,
\[
 (F-G)'=f-f=0
\]
である.前章の定理\ref{定理:導関数同じ}により,1つの区間では $F-G$ は定数になる.逆に,\,$F+C$ を微分すれば $F'=f$ に戻る.たとえば $x^2$ の原始関数は $\frac{x^3}{3}+C$ である.

なお,原始関数は1つの区間ごとに考える.たとえば $\frac1x$ の定義域は $x<0$ と $x>0$ に分かれるので,それぞれの区間で $\log|x|+C$ と書けるが,2つの区間で同じ定数を使う必要はない.また,基本定理から連続関数には原始関数が存在するが,その原始関数が既知の関数の式で簡単に表せるとは限らない.

\begin{theorem}{微積分の基本定理：原始関数で計算する}\label{定理:基本定理第二}
$f$ が $[a,b]$ で連続で,\,$F$ が $[a,b]$ で連続,\,$(a,b)$ で $F'=f$ を満たすなら,
\[
 \int_a^b f(x)\dd x=F(b)-F(a)
\]
が成り立つ.右辺を $[F(x)]_a^b$ とも書く.
\end{theorem}

証明には $A(x)=\int_a^x f(t)\dd t$ を使う.基本定理の前半から $A'=f$ なので,
\[
 (F-A)'=f-f=0
\]
である.\,$F-A$ は $[a,b]$ で連続だから,定理\ref{定理:導関数同じ}により $[a,b]$ 全体で一定となる.\,$A(a)=0$ より,その定数は $F(a)$ であり,
\[
 F(x)-A(x)=F(a)
\]
すなわち
\[
 A(x)=F(x)-F(a)
\]
を得る.ここに $x=b$ を代入すれば,求める公式となる.逆向きの積分や両端が同じ場合も,それぞれの定義から同じ式が成り立つ.

定積分は,初めから原始関数の差として定義したのではない.リーマン和の極限として定義した積分を微分すると $f$ に戻ることを示し,そこから任意の原始関数と結び付けたのである.これにより,長方形の和で表される量を,微分の逆を探す計算で求められるようになった.

証明の中で得た $A(x)=F(x)-F(a)$ は,\chapterref{章:導入}で残した問いにも答えている.\,$a=0$,\,$f(x)=x^2$,\,$F(x)=\frac{x^3}{3}$ とすれば,\,$y=x^2$ の下の面積は $A(x)=\frac{x^3}{3}-0=\frac{x^3}{3}$ となり,面積の変化率から予想した候補が確かに面積であったと分かるのだ.

同じ章の速度$v(t)=2t$にも戻ってみよう.微分して$2t$になる関数は$t^2$だから,基本定理によって
\[
 \int_0^T2t\dd t=[t^2]_0^T=T^2
\]
と求められる.長方形の和と三角形の面積から求めた値が,今度は微分の公式から得られた.一般の速度でも,原始関数が見つかれば,その端点での値の差から位置の変化を求められる.

なお,不定積分では $F+C$ と書くが,定積分では
\[
 \{F(b)+C\}-\{F(a)+C\}=F(b)-F(a)
\]
と定数が打ち消し合う.定積分の計算結果に積分定数を付ける必要はない.

\begin{bookpoint}
\textbf{POINT}\par
微分すると $f$ になる $F$ を1つ見つけ,\,$\int_a^b f(x)\dd x=F(b)-F(a)$ を計算する.
\end{bookpoint}

\begin{example}{原始関数の確認}
次の定積分を求めよ.
\[
 \text{(1)}\quad\int_0^2(3x^2-4x+1)\dd x,\qquad
 \text{(2)}\quad\int_1^{\e}\frac1x\dd x,\qquad
 \text{(3)}\quad\int_0^{\frac\pi2}\cos x\dd x.
\]
\end{example}
\paragraphbf{方針}
(1) は項ごとに,微分して $x^2,x,1$ の定数倍になるべき関数を探す.(2) は $(\log x)'=\frac1x$,(3) は $(\sin x)'=\cos x$ を逆に読む.

\paragraphbf{解答}
\begin{enumerate}[label=(\arabic*),leftmargin=*]
 \item $2$
 \item $1$
 \item $1$
\end{enumerate}

\paragraphbf{解法}
\begin{enumerate}[label=(\arabic*),leftmargin=*]
 \item 項ごとに探して $F(x)=x^3-2x^2+x$ とすれば,\,$F'(x)=3x^2-4x+1$ だから,
 \begin{align*}
  \int_0^2(3x^2-4x+1)\dd x
  &=\left[x^3-2x^2+x\right]_0^2\\
  &=(8-8+2)-0\\
  &=2
 \end{align*}
 となる.
 \item 区間 $[1,\e]$ では $x>0$ だから,\,$\log x$ を原始関数として使える.
 \begin{align*}
  \int_1^{\e}\frac1x\dd x
  &=[\log x]_1^{\e}\\
  &=1-0\\
  &=1
 \end{align*}
 \item $(\sin x)'=\cos x$ より,
 \begin{align*}
  \int_0^{\frac\pi2}\cos x\dd x
  &=[\sin x]_0^{\frac\pi2}\\
  &=1-0\\
  &=1
 \end{align*}
 となる.
\end{enumerate}

前に和から求めた $\int_1^3(2x-1)\dd x$ も,
\[
 \int_1^3(2x-1)\dd x=\left[x^2-x\right]_1^3=6
\]
と短く計算できる.2つの方法が同じ値になることは,基本定理が保証している.

\paragraphbf{平均値の定理との関係}
原始関数を使うと,\,\sectionref{節:定積分性質}で定めた関数の平均値は
\begin{align*}
 f_{\mathrm{av}}
 &=\frac1{b-a}\int_a^b f(x)\dd x\\
 &=\frac{F(b)-F(a)}{b-a}
\end{align*}
となり,原始関数 $F$ の平均変化率に等しい.したがって,積分の平均値の定理
\[
 \int_a^b f(x)\dd x=f(c)(b-a)
\]
は,\,$F'=f$ を使えば
\[
 \frac{F(b)-F(a)}{b-a}=F'(c)
\]
と書き直せる.これは $F$ に平均値の定理を使ったときの式と同じ形である.関数の平均値と平均変化率は別の量だが,原始関数を通すと,\,$f$ の平均値は $F$ の平均変化率として読めるのだ.
